\documentclass[12pt,letterpaper]{article}
\usepackage[margin=1in]{geometry}
\usepackage{enumitem}
\usepackage{amsfonts}
\usepackage[hidelinks]{hyperref}
\hypersetup{
  pdftitle={Synergistic Activities},
  pdfauthor={Juliette Bruce},
  pdfsubject={NSF DMS Algebra and Number Theory proposal}
}
\pagestyle{empty}
\setlength{\parindent}{0pt}
\setlength{\parskip}{0pt}
\setlength{\emergencystretch}{1em}
\hyphenation{VirtualResolutions Spectra Madison Explorations Zacovic}
\raggedbottom

\begin{document}
\begin{center}
  {\large\bfseries Synergistic Activities}\\[4pt]
\end{center}
\vspace{2pt}

\begin{enumerate}[label=\textbf{\arabic*.},leftmargin=0.22in,
                  labelsep=0.07in,itemsep=9pt,topsep=0pt,parsep=0pt]

\item \textbf{Open software for multigraded algebra.}
The PI co-developed \mbox{\textit{VirtualResolutions}}, an open-source package distributed with Macaulay2 for constructing and checking virtual resolutions and computing multigraded regularity on products of projective spaces. Its documented algorithms and worked examples give researchers and students practical access to methods connecting commutative algebra and algebraic geometry. Accompanied by a peer-reviewed article in the \textit{Journal of Software for Algebra and Geometry} (2020), the package translates theoretical constructions into reusable tools for investigating examples and testing conjectures.

\item \textbf{Advanced mathematical training through PAWS.}
The PI serves as lecturer for \textit{Algebraic Geometry via Computation} in the 2026 Preliminary Arizona Winter School (PAWS), a nine-week virtual program with over 300 participants, including approximately 200 students who have not yet entered Ph.D. programs. She prepared 100 pages of lecture notes on polynomial ideals,

\item \textbf{Research mentoring through MREG.}
The PI led a group of four early-stage graduate students at the 2023 Midwest Research Experience for Graduates at the University of Michigan, introducing them to research on matroids at the interface of algebraic geometry and combinatorics. The collaboration that began there led to the 2026 preprint \textit{Explorations of Matroid Complexes}, coauthored with Jacob Bucciarelli and Bailee Zacovic, and publicly available computational code and data. The project linked introductory research training with sustained collaboration and the creation of reusable computational resources.

\item \textbf{Plenary lecture at the Fields Institute.}
The PI delivered the plenary lecture \textit{Syzygies Beyond Projective Space} at the Fields Institute's \(\mathrm{Spec}(\overline{\mathbb{Q}}(2\pi i))\) conference in June 2024. Addressing a general algebraic geometry audience, she explained how homological methods over multigraded polynomial rings extend the study of syzygies beyond projective space to toric varieties. The publicly available recording makes these methods accessible to researchers and students beyond the conference, supporting the exchange of ideas between commutative algebra and algebraic geometry.

\item \textbf{National professional leadership through Spectra.}
As inaugural president of \textit{Spectra}, the PI helped develop a national professional community supporting mathematicians and students across institutions and research areas. She helped establish the annual Spectra Lavender Lecture, co-organizing its inaugural presentation at the 2022 Joint Mathematics Meetings. The lecture recognizes contributions to mathematical research, education, and the profession, giving mathematicians a recurring platform to share their work and making these contributions visible to students and the wider mathematical community.

\end{enumerate}
\end{document}
